% META: {"id":"1SPE-GEOREP-EX-034","chapitre":"1SPE-GEOMETRIE-REPEREE","type_objet":"exercice","capacites":["1SPE-GEOMETRIE-REPEREE-C4"],"capacites_codes":["C4"],"methodes":["M4"],"parcours":1,"competences":["calculer","raisonner"],"duree_min":12,"mode_creation":"ex_nihilo","sources_inspiration":[],"fichier_tex":"chapitres/1SPE-GEOMETRIE-REPEREE/exercices/1SPE-GEOREP-EX-034.tex","corrige_tex":"chapitres/1SPE-GEOMETRIE-REPEREE/corriges/1SPE-GEOREP-CO-034.tex","status":"approved","origin":"generated","status_history":["generated","audited","accepted","approved"],"release_acceptance":"RELEASE_OWNER_BATCH_ACCEPTANCE","acceptance_closure_digest":"sha256:9b3ccf9a81c5520fbb7e03b7b2d3e7bf2057a02834b6d908bf3be5f49f3b553f","acceptance_receipt_digest":"sha256:4fad86f0282e0fa4e0419cca1ad6379ae7af5a280c04a4a90880f90a1c044242"}
% BEGIN-VERIFY
% from sympy import *
% # Cercle C: centre (0,0), rayon 5. Droite D: 3x + 4y - 15 = 0
% d = abs(3*0 + 4*0 - 15) / sqrt(9+16)
% assert d == 3
% # d = 3 < 5 = r : secante
% assert 3 < 5
% END-VERIFY
\begin{exercice}{1SPE-GEOREP-EX-034}{1}{12}
Soit le cercle $\mathcal{C}$ de centre $O(0 ; 0)$ et de rayon $5$, et la droite $\mathcal{D} : 3x + 4y - 15 = 0$.

Déterminer la position relative de $\mathcal{D}$ et $\mathcal{C}$.
\end{exercice}
